\chapter{The Lambda Calculus}\label{ch:lambda}

The $\lambda$-calculus is three rules of syntax and one rule of computation, and it can compute
everything a computer can. Relearning it means relearning substitution --- the operation everyone
gets wrong once --- and the Church numerals, in which arithmetic is done by functions applied to
functions. The engine has a world for it in which every $\beta$-step is a recorded step, and a
budget, the only one in the engine, because whether a term \emph{has} a normal form is
undecidable.

\section{Terms}

\begin{definition}
  $\lambda$-terms are given by the grammar $t ::= x \mid \lambda x.\, t \mid t\, t$: variables,
  abstractions, applications. Application associates to the left and $\lambda x\, y.\, t$
  abbreviates $\lambda x.\, \lambda y.\, t$.
\end{definition}

\begin{definition}[Free variables]
  $\mathrm{FV}(x) = \{x\}$; $\mathrm{FV}(\lambda x.\, t) = \mathrm{FV}(t) \setminus \{x\}$;
  $\mathrm{FV}(t\, s) = \mathrm{FV}(t) \cup \mathrm{FV}(s)$.
\end{definition}

An occurrence of $x$ inside $\lambda x.\, t$ is \emph{bound}; the name of a bound variable is
immaterial, and $\lambda x.\, x$ and $\lambda y.\, y$ are the same function. Terms that differ
only in bound names are $\alpha$-equivalent, and the calculus is really about terms up to $\alpha$.

\section{Substitution}

$\beta$-reduction needs $t[x := s]$, the result of replacing the free occurrences of $x$ in $t$
by $s$. The naive definition --- replace every $x$ not under a $\lambda x$ --- is wrong:
$(\lambda y.\, x)[x := y]$ would become $\lambda y.\, y$, the identity, when it should be the
constant function returning $y$. The free $y$ has been \emph{captured} by the binder.

\begin{definition}[Capture-avoiding substitution]
  \begin{align*}
    x[x := s] &= s, \qquad y[x := s] = y \text{ for } y \neq x, \\
    (t\, u)[x := s] &= t[x := s]\; u[x := s], \\
    (\lambda x.\, t)[x := s] &= \lambda x.\, t, \\
    (\lambda y.\, t)[x := s] &= \lambda y.\, t[x := s] \quad\text{if } y \neq x \text{ and } y \notin \mathrm{FV}(s), \\
    (\lambda y.\, t)[x := s] &= \lambda y'.\, t[y := y'][x := s] \quad\text{otherwise, } y' \text{ fresh.}
  \end{align*}
\end{definition}

The last clause renames the binder before substituting. It is also the clause that is not
structurally recursive: it recurses on $t[y := y']$, which is not a subterm of the input. Lean
will not accept it as a total function without a separate argument that the renaming does not
change the size, and the author's earlier book on the subject marked its \lc{subst} as
\lc{partial}.

\begin{theorem}
  $\mathrm{FV}(t[x := s]) \subseteq \mathrm{FV}(s) \cup (\mathrm{FV}(t) \setminus \{x\})$.
\end{theorem}
\begin{proof}
  Induction on $t$, clause by clause. The only case with content is the abstraction: a variable
  free in $\lambda y.\, t[x := s]$ is free in $t[x := s]$ and is not $y$, so by induction it is in
  $\mathrm{FV}(s)$ or in $\mathrm{FV}(t) \setminus \{x\}$, and in the second case it is in
  $\mathrm{FV}(\lambda y.\, t) \setminus \{x\}$ since it is not $y$.
\end{proof}

\section{Reduction}

\begin{definition}[$\beta$]
  $(\lambda x.\, t)\, s \to t[x := s]$, applied anywhere inside a term. A term with no
  $\beta$-redex is in normal form.
\end{definition}

Two facts, stated without proof, organize everything about reduction. \emph{Church--Rosser}: if
$t \to^* u$ and $t \to^* v$ then some $w$ has $u \to^* w$ and $v \to^* w$; so a normal form, if it
exists, is unique. \emph{Standardization}: if $t$ has a normal form, \emph{normal-order}
reduction --- always contract the leftmost-outermost redex --- reaches it. The second is why the
engine reduces in normal order: any other strategy might loop on a term that has a normal form.

\begin{theorem}
  $\Omega = (\lambda x.\, x\, x)(\lambda x.\, x\, x)$ has no normal form.
\end{theorem}
\begin{proof}
  Its only redex is the whole term, which reduces to itself.
\end{proof}

And the general question is hopeless: since the calculus can express every computable function,
``does $t$ have a normal form'' is undecidable, by reduction from the halting problem. An engine
that promises an answer for every term is lying; the honest engine has a budget and says so.

\section{Church numerals}

\begin{definition}
  $\overline{n} = \lambda f\, x.\, f^n\, x$, where $f^n\, x$ is $f$ applied $n$ times to $x$.
  $\mathrm{succ} = \lambda n\, f\, x.\, f\,(n\, f\, x)$; $\mathrm{add} = \lambda m\, n\, f\, x.\,
  m\, f\, (n\, f\, x)$; $\mathrm{mul} = \lambda m\, n\, f.\, m\, (n\, f)$.
\end{definition}

A numeral is an iterator: $\overline{n}\, f\, x$ applies $f$ to $x$ $n$ times. Addition iterates
$f$ $n$ times and then $m$ more; multiplication iterates ``iterate $f$ $n$ times'' $m$ times.

\begin{exbox}[$\mathrm{add}\, \overline 2\, \overline 3$]
  $\mathrm{add}\, \overline 2\, \overline 3 \to^2 \lambda f\, x.\, \overline 2\, f\, (\overline 3\, f\, x)
  \to^2 \lambda f\, x.\, \overline 2\, f\, (f (f (f x))) \to^2 \lambda f\, x.\, f (f (f (f (f x)))) =
  \overline 5$. Seven $\beta$-steps in all, counting the unfolding of the numerals' own
  applications.
\end{exbox}

Booleans are selectors, $\mathrm{true} = \lambda t\, f.\, t$ and $\mathrm{false} = \lambda t\, f.\,
f$, and $\mathrm{if}\, b\, u\, v = b\, u\, v$.

\section{De Bruijn indices}

Names are for people. For a machine, a bound variable can be written as the number of binders
between its occurrence and its own: $\lambda x.\, \lambda y.\, x$ is $\lambda\, \lambda\, 1$ and
$\lambda x.\, x\, (\lambda y.\, x\, y)$ is $\lambda\, 0\, (\lambda\, 1\, 0)$. $\alpha$-equivalent
terms become \emph{identical}, and substitution becomes arithmetic on indices. The engine
computes this view beside every step so the notebook can toggle it.

\section{The engine}

\begin{minted}{lean}
inductive Term where
  | var : String → Term
  | lam : String → Term → Term
  | app : Term → Term → Term

def freshen (clash : List String) (ren : List (String × String)) : Term → Term
  | .var w => .var ((ren.lookup w).getD w)
  | .app a b => .app (freshen clash ren a) (freshen clash ren b)
  | .lam y e =>
    let y' := if y ∈ clash then freshVar (clash ++ freeVars e ++ ren.map (·.2)) y else y
    .lam y' (freshen clash ((y, y') :: ren) e)

def substRaw (x : String) (s : Term) : Term → Term
  | .var y => if y == x then s else .var y
  | .app a b => .app (substRaw x s a) (substRaw x s b)
  | .lam y e => if y == x then .lam y e else .lam y (substRaw x s e)

def subst (x : String) (s : Term) (e : Term) : Term × Bool :=
  let clash := (freeVars s).filter (· != x)
  let e' := freshen clash [] e
  (substRaw x s e', e' != e)
\end{minted}

This is the definition of the second section rearranged to be structurally recursive: first one
pass renames every binder that could capture (\lc{freshen}, carrying the renaming down as a
parameter), then one pass substitutes without renaming (\lc{substRaw}). Each pass is structural,
so the whole is total, and it computes the same term as the textbook's clause. The Boolean says
whether any renaming happened, so the notebook can label the step an $\alpha$-then-$\beta$.

\begin{minted}{lean}
def betaStep : Term → Option (Term × Bool)
  | .app (.lam x body) arg => some (subst x arg body)
  | .app a b =>
    match betaStep a with
    | some (a', r) => some (.app a' b, r)
    | none => (betaStep b).map fun (b', r) => (.app a b', r)
  | .lam x e => (betaStep e).map fun (e', r) => (.lam x e', r)
  | .var _ => none

def maxSteps : Nat := 1000
\end{minted}

\lc{betaStep} is normal order: the leftmost-outermost redex, found by trying the head of an
application first. \lc{reduce} iterates it up to \lc{maxSteps} and reports the term it reached if
no normal form came; $\Omega$ is refused after $1000$ steps, still $\Omega$.

\section{What is proved}

\begin{minted}{lean}
theorem substRaw_freeVars (x : String) (s : Term) : ∀ (e : Term) (z : String),
    z ∈ freeVars (substRaw x s e) → z ∈ freeVars s ∨ (z ∈ freeVars e ∧ z ≠ x)

theorem readChurch_church (n : Nat) : readChurch (church n) = some n
\end{minted}

The first is the theorem of the second section, for the capture-free pass. The second says the
reader that labels a normal form ``the Church numeral $n$'' is right on the numerals the engine
builds --- a small theorem, but the one that makes the label beside \texttt{add 2 3} a claim
rather than a guess.

Not proved: that \lc{freshen} produces an $\alpha$-equivalent term, hence that a $\beta$-step
under a renaming is a $\beta$-step of the original; and that the de Bruijn view commutes with
$\beta$. The $\beta$-steps are reported \status{unverified} for that reason, and the unfolding of
a definition (\rn{lambda.delta}) \status{verified}, because it is a lookup with nothing to prove.
The missing metatheory is standard --- it is the content of the first chapters of any book on
the subject --- and is the natural next thing to add.

\begin{trybox}
\begin{minted}{text}
add 2 3
mul 2 3
if true 1 0
(λx y. x) a b
(λx. x x) (λx. x x)
\end{minted}
  Type \texttt{\textbackslash lam} and Tab for the $\lambda$. Toggle de Bruijn indices in the
  View menu and watch the names disappear. The last is $\Omega$, refused.
\end{trybox}

\section{Exercises}

\begin{exercisebox}[Capture]
  Compute $(\lambda y.\, x\, y)[x := y]$ by the naive definition and by the capture-avoiding
  one, and apply both results to a variable $z$ to see the difference.
\end{exercisebox}

\begin{exercisebox}[Multiplication]
  Reduce $\mathrm{mul}\, \overline 2\, \overline 3$ by hand in normal order and count the
  steps. Then find the reduction that is shortest, and explain why it is not normal order.
\end{exercisebox}

\begin{exercisebox}[Two passes equal one]
  Prove, on paper, that \lc{substRaw x s (freshen clash [] e)} equals the textbook's
  $e[x := s]$ up to $\alpha$, when \lc{clash} is $\mathrm{FV}(s) \setminus \{x\}$. Where does the
  proof use that the fresh names avoid $\mathrm{FV}(e)$?
\end{exercisebox}

\begin{exercisebox}[Indices]
  Write $\mathrm{add}$ and $\overline 3$ with de Bruijn indices, and write the rule for
  substituting under a binder as a shift of indices.
\end{exercisebox}
